\subsection{动量有限差分与指数最优性}
\label{app:hamilton-finite-difference}

下面证明前一小节使用的动量差分估计。全程保留原来的非 Lipschitz 梯度，不使用平滑正则化。
记
\[
 g_1(t)=\operatorname{sgn}(t)\sqrt{|t|},\qquad
 h(p)=\begin{cases}p/\sqrt{\|p\|},&p\ne0,\\0,&p=0,\end{cases}
 \qquad \beta=\frac{\sqrt2}{4}.
\]
其中 $g_1$ 连续且严格递增。近端输入 $(a,b)$ 与输出 $(x,p)$
满足
\begin{equation}
\label{eq:hamilton-app-prox}
 x=a+h(p),\qquad b=g(x)+100p,
 \qquad g(x)=(g_1(x_1),0).
\end{equation}

\begin{lemma}[动量球上的有限差分]
\label{lem:hamilton-secant}
设 $p=(q,r)$、$\widetilde p=(\widetilde q,\widetilde r)$
均满足 $\|p\|,\|\widetilde p\|\leq R$，其中 $R>0$。
记 $\Delta q=q-\widetilde q$、$\Delta r=r-\widetilde r$。
若 $\Delta q\geq\beta|\Delta r|$，则
\begin{equation}
\label{eq:hamilton-secant}
 h_1(p)-h_1(\widetilde p)
 \geq\frac{\Delta q-\beta|\Delta r|}{2\sqrt R}.
\end{equation}
\end{lemma}

\begin{proof}
令 $p_\tau=\widetilde p+\tau(p-\widetilde p)$，$0\leq\tau\leq1$。
球的凸性给出 $\|p_\tau\|\leq R$。
在 $p_\tau\ne0$ 的点，直接对
$h_1(q,r)=q(q^2+r^2)^{-1/4}$ 求导，得到
\[
 A(q,r):=\frac{\partial h_1}{\partial q}
 =\frac{q^2+2r^2}{2(q^2+r^2)^{5/4}}
 \geq\frac1{2\sqrt R},\qquad
 B(q,r):=\frac{\partial h_1}{\partial r}
 =-\frac{qr}{2(q^2+r^2)^{5/4}}.
\]
又因为 $q^2+2r^2\geq2\sqrt2|qr|$，有
$|B|\leq\beta A$。故在上述可微点，
\[
 \frac{d}{d\tau}h_1(p_\tau)
 =A(p_\tau)\Delta q+B(p_\tau)\Delta r
 \geq\frac{\Delta q-\beta|\Delta r|}{2\sqrt R}.
\]
若线段不经过原点，普通微积分基本定理即给出
\eqref{eq:hamilton-secant}。

若 $p=\widetilde p=0$，两端均为零。否则线段至多在一个
参数 $\tau_0$ 处经过原点。写 $v=p-\widetilde p\ne0$，则
$p_\tau=(\tau-\tau_0)v$，并且有精确公式
\[
 h_1(p_\tau)
 =\operatorname{sgn}(\tau-\tau_0)
   \sqrt{|\tau-\tau_0|}\,\frac{v_1}{\sqrt{\|v\|}}.
\]
因此它在 $\tau_0$ 连续。分别在原点两侧去掉长度为
$\varepsilon$ 的小区间，积分已经证明的导数下界，再由此精确
公式令 $\varepsilon\downarrow0$，便得到
\eqref{eq:hamilton-secant}。$\tau_0=0$ 或 $1$ 时只需去掉
相应端点的小区间。这一过程只是对原函数的差分求界，不涉及
正则化函数或平滑极限。
\end{proof}

\begin{lemma}[近端输出动量的两点估计]
\label{lem:hamilton-momentum-difference}
设 $(x,p)$、$(\widetilde x,\widetilde p)$ 分别对应近端输入
$(a,b)$、$(\widetilde a,\widetilde b)$，且
$\|p\|,\|\widetilde p\|\leq R$。则
\begin{align}
\label{eq:hamilton-momentum-first-coordinate}
 |p_1-\widetilde p_1|
 &\leq\max\left\{
       \frac{|b_1-\widetilde b_1|}{100},\,
       2\sqrt R|a_1-\widetilde a_1|
       +\frac{\beta}{100}|b_2-\widetilde b_2|
      \right\},\\
\label{eq:hamilton-momentum-full-difference}
 \|p-\widetilde p\|
 &\leq\left(4R+\frac{2+\beta^2}{10000}\right)^{1/2}
       \|(a,b)-(\widetilde a,\widetilde b)\|.
\end{align}
特别地，在 $R=10^{-6}$ 时，右侧常数小于 $0.015$，也小于 $0.02$。
\end{lemma}

\begin{proof}
由 \eqref{eq:hamilton-app-prox}，
$p_2-\widetilde p_2=(b_2-\widetilde b_2)/100$。
以下先设 $\Delta q:=p_1-\widetilde p_1>0$。
若 $\Delta q$ 严格大于
\eqref{eq:hamilton-momentum-first-coordinate} 右端的两个量，
则引理~\ref{lem:hamilton-secant} 给出
\[
 h_1(p)-h_1(\widetilde p)>|a_1-\widetilde a_1|,
 \qquad x_1-\widetilde x_1>0.
\]
由 $g_1$ 的严格递增性，
$g_1(x_1)-g_1(\widetilde x_1)>0$，从而
\[
 b_1-\widetilde b_1
 =100\Delta q+g_1(x_1)-g_1(\widetilde x_1)
 >100\Delta q>|b_1-\widetilde b_1|,
\]
矛盾。交换两点可处理 $\Delta q<0$，$\Delta q=0$ 时左端为零。
上述比较只用 $g_1$ 的单调性，所以同样适用于
$x_1=0$、$\widetilde x_1=0$，以及两点跨越零点的情形。

将最大值放大为非负项之和，有
\[
 |p_1-\widetilde p_1|
 \leq2\sqrt R|a_1-\widetilde a_1|
       +\frac{|b_1-\widetilde b_1|}{100}
       +\frac{\beta|b_2-\widetilde b_2|}{100}.
\]
对此应用 Cauchy--Schwarz 不等式，并加上第二动量坐标的平方，
得到 \eqref{eq:hamilton-momentum-full-difference}。
在 $R=10^{-6}$ 时，该常数的平方为 $0.0002165<0.015^2$。
\end{proof}

\begin{lemma}[动能梯度的全局 Hölder 界]
\label{lem:hamilton-h-holder}
对任意 $p,\widetilde p\in\mathbb R^2$，有
\begin{equation}
\label{eq:hamilton-h-holder}
 \|h(p)-h(\widetilde p)\|^2\leq2\|p-\widetilde p\|.
\end{equation}
\end{lemma}

\begin{proof}
令 $u=h(p)$、$v=h(\widetilde p)$，则
$p=\|u\|u$、$\widetilde p=\|v\|v$，包括零向量在内。
记 $\alpha=\|u\|$、$\eta=\|v\|$。直接展开可得
\[
 \langle p-\widetilde p,u-v\rangle
 =\frac{\alpha+\eta}{2}
   \bigl(\|u-v\|^2+(\alpha-\eta)^2\bigr)
 \geq\frac12\|u-v\|^3,
\]
其中使用了 $\alpha+\eta\geq\|u-v\|$。
再用 Cauchy--Schwarz 不等式并在 $u\ne v$ 时约去
$\|u-v\|$，即得 \eqref{eq:hamilton-h-holder}；$u=v$ 时两端均非负且左端为零。
\end{proof}

\begin{proposition}[指数的最优性与上界的保守性]
\label{prop:hamilton-sharp-exponents}
对于命题~\ref{prop:rank-deficient-hamilton-rleb} 的完整图块，
幂次 RL 的最大指数以及完整预解式 $J_F$ 的最大局部 H\"older 指数均为 $1/2$。
真残差误差界中幂次 $2$ 也不能提升为任何大于 $2$ 的幂次。
距离收缩与尾项估计中的数值常数和几何因子仅作为上界使用。
\end{proposition}

\begin{proof}
取 $w_t=((0,0),(0,t))$，$0<t\le R$，以及 $w_0=0$。
方程~\eqref{eq:hamilton-app-prox} 给出
\[
 y_t:=J_F(w_t)=((0,\sqrt{t/100}),(0,t/100)),\qquad J_F(w_0)=0.
\]
两输入均在 $U_R$，且 $\|w_t-w_0\|=t$。
对这一图点与零点，直接计算得
\[
 E_t:=\|y_t\|^2+\|F(y_t)\|^2
 =\frac{t}{50}+\frac{4901t^2}{5000},\qquad
 2E_t-t^2=\frac{t}{25}+\left(\frac{49t}{50}\right)^2.
\]
若该图块满足指数 $\gamma>1/2$、有限常数 $L'$ 的幂次 RL 条件，
引理~\ref{lem:RL-J} 要求
\[
 (L')^2t^{2\gamma}\ge2E_t-t^2\ge\frac{t}{25},
\]
这与 $t\downarrow0$ 矛盾。
结合已证明的 $1/2$ 阶估计，RL 最大指数恰为 $1/2$。
同一计算给出该阶任何统一 RL 常数至少为 $1/5$，
而上文所取 $L=1/2$ 并不声称最优。
另一方面，$\|J_F(w_t)-J_F(w_0)\|\ge\sqrt t/10$
排除完整预解式的任何高于 $1/2$ 的局部 H\"older 指数；
平方能量上界则给出其 $1/2$ 阶界。

再取 $y_t=((t^2,0),0)$，$t>0$。
其真残差为 $r_F(y_t)=t$，距离为 $d(y_t,S)=t^2$。
对应近端输入为 $y_t+F(y_t)=((t^2,0),(t,0))$，其到 $S$
的距离为 $\sqrt{t^4+t^2}$，故充分小的 $t$ 给出图块中的输出点。
于是任何 $d(y,S)\leq K r_F(y)^{2+\varepsilon}$，
$\varepsilon>0$，都将要求 $1\leq Kt^\varepsilon$，不可能成立。

最后，对上述初值 $z^0=w_t$，可直接求出整条轨道：
\[
 z^j=\left(
      \left(0,\frac{\sqrt t}{9}(1-10^{-j})\right),
      (0,100^{-j}t)\right),\qquad
 z^\infty=((0,\sqrt t/9),0).
\]
因此实际距离和尾项分别为
\[
 d(z^j,S)=100^{-j}t,\qquad
 \|z^j-z^\infty\|
 =\left(\frac{t}{81}100^{-j}+t^2 10000^{-j}\right)^{1/2}.
\]
这条轨道满足正文的 $4^{-j}$ 距离界与 $2^{-j}$ 阶尾界，
但其实际收敛更快，因而不能把正文上界当作轨道的精确渐近式。
\end{proof}
